\documentclass[10pt]{article}
\usepackage[margin=1in]{geometry}
\usepackage[T1]{fontenc}
\usepackage{lmodern,microtype,amsmath,amssymb,amsthm}
\theoremstyle{definition}
\newtheorem{assumption}{Assumption}
\begin{document}
\section*{Constrained Monotone Operators with Gap Noise}
We consider a variational inequality for a continuous monotone operator $F:E\to E$ on a convex set $X$ in a finite-dimensional real Hilbert space $E$. The operator satisfies a generalized Lipschitz bound, and a given point $x^0\in X$ lies near a solution. Two oracles return independent-sample unbiased estimates of $F$ on $X$, with variance controlled by an operator gap, and exact projections onto $X$. For $\varepsilon>0$ and $\delta\in(0,1)$, the goal is a measurable output $\bar{x}\in X$ admitting $u\in E$ with $\langle u,w-\bar{x}\rangle\leq0$ for every $w\in X$ and $\|F(\bar{x})+u\|\leq\varepsilon$, with probability at least $1-\delta$.

The underlying variational inequality is to find $x^\star\in X$ such that
\[
\langle F(x^\star),w-x^\star\rangle\geq0
\qquad\text{for every }w\in X.
\]
The norm is induced by the inner product on $E$. The parameters satisfy
\[
L_0>0,\qquad L_1\geq0,\qquad R>0,\qquad
G,\sigma,v\geq0.
\]
The algorithm receives $E$, a measurable sample space $\Omega$, $x^0$, and all parameters $L_0,L_1,R,G,\sigma,v,\varepsilon,\delta$; it accesses $F$ and $X$ through the two oracles.

Our assumptions:
\begin{assumption}[Feasibility and reference solution]
The set $X\subseteq E$ is convex and contains $x^0$. There exists a solution $x^\star$ of the variational inequality satisfying
\[
\|x^0-x^\star\|\leq R,
\qquad \|F(x^\star)\|\leq G.
\]
An exact projection map $\Pi_X:E\to E$ is available and satisfies, for every $z\in E$,
\[
\begin{aligned}
\Pi_X(z)&\in X,\\
\langle z-\Pi_X(z),w-\Pi_X(z)\rangle&\leq0
\qquad\text{for every }w\in X.
\end{aligned}
\]
\end{assumption}

\begin{assumption}[Monotonicity and generalized Lipschitz bound]
The operator $F$ is continuous on $E$, and for all $x,y\in E$,
\[
\langle F(x)-F(y),x-y\rangle\geq0.
\]
For every $x,y\in E$ and $M\in\mathbb{R}$, the segment bound
\[
\|F(x+\theta(y-x))\|\leq M
\qquad\text{for every }\theta\in[0,1]
\]
implies
\[
\|F(x)-F(y)\|\leq (L_0+L_1M)\|x-y\|.
\]
\end{assumption}

\begin{assumption}[Gap-dependent stochastic noise]
There is a probability measure $\nu$ on $\Omega$ and a jointly measurable map $g:E\times\Omega\to E$. For every $x\in X$, the random vector $g(x,Z)$ is Bochner integrable and $\|g(x,Z)-F(x)\|^2$ is integrable for $Z\sim\nu$, with
\[
\begin{aligned}
\mathbb{E}_{Z\sim\nu}[g(x,Z)]&=F(x),\\
\mathbb{E}_{Z\sim\nu}\bigl[\|g(x,Z)-F(x)\|^2\bigr]
&\leq \sigma^2+v\langle F(x)-F(x^\star),x-x^\star\rangle.
\end{aligned}
\]
The reference solution $x^\star$ is the same as above. The oracle samples $(Z_k)_{k\geq0}$ are measurable, mutually independent, and each has law $\nu$. This sequence is independent of the algorithm's randomness.
\end{assumption}

Both oracles accept queries in $E$. At the $k$th oracle call, a stochastic query at $x$ returns $g(x,Z_k)$, whereas a projection query at $z$ returns $\Pi_X(z)$. The stochastic moment conditions apply to queries in $X$. Stochastic-operator calls and projection calls are counted separately.

Define the normal cone at $x\in X$ by
\[
N_X(x)=\{u\in E:\langle u,w-x\rangle\leq0
\text{ for every }w\in X\}.
\]
After the allotted calls, the output $\bar{x}$ must be measurable and satisfy
\[
P\!\left(\neg\left[\bar{x}\in X\ \text{and}\
\exists u\in N_X(\bar{x}):\|F(\bar{x})+u\|\leq\varepsilon\right]\right)
\leq\delta,
\]
where $P$ is the joint law of the oracle samples and the independent algorithm randomness.
\end{document}
